
\section*{Introduction}

In this chapter we will try to develop some of the derived theory needed to handle condensed phenomena later on, \eg in the definition of analytic rings in definition \ref{definition: analytic rings}. For the derived theory of sheaves of modules on a scheme (a setting with plenty of injectives), we generically refer to the \cite[The Stacks Project]{stacks-project} which covers all needed formalism extensively.

Why does the author write \emphquote{try to}? In the main reference for the material covered in this thesis -- the lecture notes \cite{condenseddotpdf} -- Clausen and Scholze are quite lenient in their use of derived algebra. They seldomly show that objects (like derived functors) exist and why they satisfy certain properties. While they occasionally give short hints, \eg that one can construct (left) derived tensor products and (right) derived internal $\Hom$s by \quote{[u]sing projective resolutions} \cite[p. 13 (iii)]{condenseddotpdf}, these hints are certainly not comprehensive. It is to note that Clausen and Scholze (in their lectures) seem to mainly work in the language of $\infty$-categories, implicit in the first chapters but explicit in \cite[Lectures IX, X \& XI]{condenseddotpdf} where the derived categories of interest are explicitly considered as \emph{stable $\infty$-categories} (which then allows one to glue them). The author is sure that all claims made by Clausen and Scholze (either implicitly or explicitly) are \quote{obvious enough} when one works (and is used to) $\infty$-categorical language -- this however does not apply to the author. 

As there are (too) many things to say about $\infty$-land for one thesis, the author decided on a more straightforward approach. We will try to introduce all necessary objects and prove their relevant properties on the fly, using more \quote{traditional} theory. This however did not work in all cases! This will be indicated (obviously) at each such place. We will however still use these results later on, even if we do not have a proof. Then with the peace in mind that Clausen and Scholze can provide a proof in $\infty$-world.

%The main reason why the author was not able to prove all results is relatively \quote{uniform}. Recall that various categories of condensed objects (in particular condensed modules over some condensed ring) do not have enough injective objects (see remark \ref{remark: not enough injectives}). Instead they have \emph{enough projectives}. However the projectives do not behave as nicely as modules over a ring: Tensor products of projectives do not need to be projective (in fact are almost never, see \cite[Prop. 3.7 \& Rem. 3.8]{complexdotpdf}). A sample consequence is that if there is an additional closed(!) symmetric monoidal structure present (as with condensed abelian groups) that these projectives are almost never \emph{internally projective}, \ie the internal $\intHom(P, -)$ of some projective $P$ is almost never exact. This makes \eg proving the existence of enriched (to internal $\intRHom$s) unnecessarily difficult.

\section{Prerequisites}

We assume the reader familiar with the basic notions and definitions of (bounded -- either to the left, to the right, or in both directions) derived algebra. This includes
\begin{itemize}
	\item 
	the definition of \emph{co-chain complexes} in some abelian category $\category{A}$ (which we will simply call \emph{chain complexes} or \emph{complexes} as we will solely work with such) and the definition of \emph{chain maps} as their morphisms, constituting a category $\complexes{\category{A}}$ (see \cite[\href{https://stacks.math.columbia.edu/tag/010V}{Tag 010V}]{stacks-project} for both),
	
	\item 
	the notion of \emph{homotopy of chain maps} (see \cite[\href{https://stacks.math.columbia.edu/tag/0112}{Tag 0112}]{stacks-project}) which leads to the formation of the \emph{homotopy category of complexes} $\homotopyCategoryOfComplexes{\category{A}}$ (see \cite[\href{https://stacks.math.columbia.edu/tag/05RN}{Tag 05RN}]{stacks-project}),
	
	\item
	the definition of \emph{triangulated categories} (see \cite[\href{https://stacks.math.columbia.edu/tag/05QK}{Tag 05QK}]{stacks-project}) with example $\homotopyCategoryOfComplexes{\category{A}}$ (see \cite[\href{https://stacks.math.columbia.edu/tag/014S}{Tag 014S}]{stacks-project}), as well as the notion of \emph{exact/triangulated functor} (see \cite[\href{https://stacks.math.columbia.edu/tag/014V}{Tag 014V}]{stacks-project}) between such categories,
	
	\item
	the notion of \emph{acyclicity}, the definition of \emph{quasi-isomorphisms} (see \cite[\href{https://stacks.math.columbia.edu/tag/010Z}{Tag 010Z}]{stacks-project} for both) and the general theory of \emph{localization of triangulated categories} at a multiplicative system compatible with the triangulated structure (see \cite[\href{https://stacks.math.columbia.edu/tag/05R1}{Tag 05R1}]{stacks-project}), leading to the definition of the \emph{derived category} $\derived{\category{A}}$ as the localization of the triangulated category $\homotopyCategoryOfComplexes{\category{A}}$ at the multiplicative system $\Sigma$ of all quasi-isomorphisms of complexes in $\category{A}$ (see \cite[\href{https://stacks.math.columbia.edu/tag/05RU}{Tag 05RU}]{stacks-project}),
	
	\item
	the definitions of \emph{left \& right resolutions} and the notions of \emph{having enough projectives} and \emph{having enough injectives} (see \cite[\href{https://stacks.math.columbia.edu/tag/0643}{Tag 0643}]{stacks-project} for the projective case, the injective case is formally dual),
	
	\item
	the definition of \emph{left and right derived functors} as the natural extension of an exact functor between triangulated categories to a localization of its domain (see \cite[\href{https://stacks.math.columbia.edu/tag/05S7}{Tag 05S7}]{stacks-project}) and how well chosen left and right resolutions allow to \emph{compute} (see \cite[\href{https://stacks.math.columbia.edu/tag/05SX}{Tag 05SX}]{stacks-project} and \cite[\href{https://stacks.math.columbia.edu/tag/06XN}{Tag 06XN}]{stacks-project}) the values of derived functors, \eg by left or right resolutions of projectives and injectives respectively.
\end{itemize}

\section{\texorpdfstring{$K$}{K}-projectivity}
\label{section: K-projectivity}

In this thesis our main application of the derived theory are of course categories of condensed objects, \eg condensed abelian group or condensed modules over a condensed ring. As already mentioned these usually do not admit enough (or rather any non-zero) injectives (see remark \ref{remark: not enough injectives}). Consequentially, we will restrict our attention to projectives (of which there are plenty since all these categories are generated by compact projectives). It is to note that the main reference for this section, \cite[The Stacks Project]{stacks-project}, works with a very restrictive class of abelian categories, see \cite[\href{https://stacks.math.columbia.edu/tag/09PA}{Tag 09PA}]{stacks-project} -- into which our categories of desire do not fit. As mentioned in the tag, if one wishes to work with a \emph{large} abelian category, one must ensure that the derived category, as a localization of the homotopy category of complexes $\homotopyCategoryOfComplexes{\category{A}}$, is locally small. This will be the content of corollary \ref{corollary: local smallness of the derived category}. Furthermore, in this thesis we will predominantly work with \emph{unbounded} derived categories. While in right bounded case its is enough to work with right bounded projective resolutions, in the unbounded case one has to be more careful. This leads to the notion of \emph{$K$-projectivity}, originally introduced by Spaltenstein in \cite{spaltenstein}. The author wants to thank \href{https://zmao-math.github.io/}{Zhouhang Mao} for introducing him to Spaltenstein's work. We will mainly follow The Stacks Project \cite{stacks-project} and only reference Spaltenstein's \cite{spaltenstein} if necessary.
\begin{definition}[K-projectivity, {\cite[\href{https://stacks.math.columbia.edu/tag/070H}{Tag 070H}]{stacks-project}}]
	
	Let $\category{A}$ be an abelian category. A complex $P^\bullet$ in $\category{A}$ is called \emph{$K$-projective} if for any acyclic complex $N^\bullet$ in $\category{A}$ we have that $\Hom_\homotopyCategoryOfComplexes{\category{A}}(P^\bullet, N^\bullet)\cong 0$.
\end{definition}

We obtain the following very useful characterization of K-projectivity.
\begin{lemma}[Morphisms out of a $K$-projective complex, {\cite[\href{https://stacks.math.columbia.edu/tag/070I}{Tag 070I}]{stacks-project}}]
	\label{lemma: morphism out of a K-projective complex}
	
	Let $\category{A}$ be an abelian category. A complex $P^\bullet$ in $\category{A}$ is $K$-projective if and only if for all complexes $N^\bullet$ in $\category{A}$ the natural map
	$$\Hom_\homotopyCategoryOfComplexes{\category{A}}(P^\bullet, N^\bullet) \to \Hom_\derived{\category{A}}(P^\bullet, N^\bullet)$$
	is a bijection.
\end{lemma}

An immediate consequence is that if the abelian category $\category{A}$ \quote{has enough $K$-projective resolutions} then the derived category $\derived{\category{A}}$ is locally small.
\begin{corollary}[Local smallness of the derived category]
	\label{corollary: local smallness of the derived category}
	
	Let $\category{A}$ be an abelian complex such that every complex admits a left resolution by a $K$-projective complex. Then $\derived{\category{A}}$ is locally small.
\end{corollary}
\begin{proof}
	Let $M^\bullet$ and $N^\bullet$ be arbitrary complexes. By assumption there is a left resolution $P^\bullet \to M^\bullet$ of $M^\bullet$ by a $K$-projective complex $P^\bullet$. Then by lemma \ref{lemma: morphism out of a K-projective complex} we obtain that
	$$\Hom_\derived{\category{A}}(M^\bullet, N^\bullet) \isorightarrow \Hom_\derived{\category{A}}(P^\bullet, N^\bullet) \isoleftarrow \Hom_\homotopyCategoryOfComplexes{\category{A}}(P^\bullet, N^\bullet)$$
	and thus that $\Hom_\derived{\category{A}}(M^\bullet, N^\bullet)$ is small.
\end{proof}

Furthermore, $K$-projective complexes \emph{compute} left derived functors.
\begin{lemma}[Existence of left derived functors, {\cite[\href{https://stacks.math.columbia.edu/tag/070K}{Tag 070K}]{stacks-project}}]
	\label{lemma: existence of left derived functors}
	
	Let $\category{A}$ be an abelian category such that every complex admits a left resolution by a $K$-projective complex and let $F\colon \homotopyCategoryOfComplexes{\category{A}} \to \category{D}'$ be an exact functor of triangulated categories. Then $\LeftD F\colon \derived{\category{A}} \to \category{D}'$ exists and for any $K$-projective complex $P^\bullet$ we obtain that the natural morphism $\LeftD F(P^\bullet) \to F(P^\bullet)$ is a quasi-isomorphism.
\end{lemma}

Spaltenstein now gives the following proposition that ensures the existence of \emph{enough $K$-projective complexes}. Recall that a \emph{sequential} (co)limit is a (co)limit over a diagram indexed by the partially ordered set $(\N, \le)$.
\begin{proposition}[Enough projectives on steroids]
	\label{proposition: enough projectives on steroids}
	
	If $\category{A}$ is an abelian category with enough projectives such that filtered colimits are exact, then every complex in $\category{A}$ admits a left resolution by a $K$-projective complex that is a sequential colimit of bounded to the right, term-wise projective complexes.
\end{proposition}
\begin{proof}
	By \cite[Example. 3.2 (a)]{spaltenstein} every bounded to the right complex in $\category{A}$ admits a left resolution by projectives and any such complex of projectives is $K$-projective. Denote by $\mathfrak P$ the class of bounded to the right complexes of projectives. Denote further the closure of $\mathfrak{P}$ under \emph{special direct limits} (see \cite[Def. 2.6]{spaltenstein}) by $\mathfrak P^\rightarrow$.
	
	Now since filtered colimits are exact, theorem \cite[Thm. 3.4]{spaltenstein} and its corollary \cite[Cor. 3.5]{spaltenstein} assert that every complex of $\category{A}$ admits a left $\mathfrak P^\rightarrow$-resolution of the desired form and that any such complex is $K$-projective.
\end{proof}

A useful result is the following.
\begin{lemma}[Preservation of $K$-projectivity, {\cite[\href{https://stacks.math.columbia.edu/tag/08BJ}{Tag 08BJ}]{stacks-project}}]
	\label{lemma: preservation of K-projectivity}
	
	Let $F\colon \category{A} \to \category{B}$ and $G\colon \category{B} \to \category{A}$ be a pair of additive functors such that $F$ is exact and left adjoint to $G$. Then $G$ preserves $K$-projectivity, \ie for any $K$-projective complex $P^\bullet$ the complex $G(P^\bullet)$ is $K$-projective as well.
\end{lemma}

%This has a direct application to the morphism of derived functors underlying Grothendieck's spectral sequence.
%\begin{lemma}[Grothendieck's isomorphism]
%	\label{lemma: grothendieck's isomorphism}	
%	
%	Let $F\colon \category{A}\to \category{B}$ and $G\colon \category{B}\to\category{C}$ be right exact functors such that over $\category{A}$ and $\category{B}$ every complex admits a $K$-projective left resolution.
%	
%	The following are equivalent
%	\begin{enumerate}
%		\item
%		$F(P^\bullet)$ computes $\LeftD G(F(P^\bullet))$ (\ie $\LeftD G(F(P^\bullet))\to G(F(P^\bullet))$ is an isomorphism) for any $K$-projective complex $P^\bullet$ in $\category{A}$.
%		
%		\item 
%		The canonical $\LeftD G \circ \LeftD F \to \LeftD(G\circ F)$ is an isomorphism of (unbounded) derived functors.
%	\end{enumerate}
%\end{lemma}
%\begin{proof}
%	This is essentially \cite[\href{https://stacks.math.columbia.edu/tag/015M}{Tag 015M}]{stacks-project}.
%	\begin{todo}
%		Check again.
%	\end{todo}
%\end{proof}

%\begin{corollary}
%	In the setting of lemma \ref{lemma: grothendieck's isomorphism} if $F$ admits an exact left adjoint, then $F$ preserves $K$-projective complexes by lemma \ref{lemma: preservation of K-projectivity}. Then every such $F(P^\bullet)$ is $K$-projective as well and thus computes $\LeftD G(F(P^\bullet))$, implying that $\LeftD G\circ \LeftD F\to \LeftD(G\circ F)$ is an isomorphism.
%\end{corollary}

%\begin{lemma}
%	Let $F\colon \category{A}\to \category{B}$ be an additive functor of abelian categories. The natural extensions of $F$ to functors $\homotopyCategoryOfComplexes{\category{A}}\to\homotopyCategoryOfComplexes{\category{B}}$ and $\homotopyCategoryOfComplexes{\category{A}}\to\derived{\category{B}}$ are exact functors of triangulated categories.
%\end{lemma}
%\begin{proof}
%	This follows from \cite[\href{https://stacks.math.columbia.edu/tag/014X}{Tag 014X}]{stacks-project} as the localization functor $\homotopyCategoryOfComplexes{\category{B}}\to\derived{\category{B}}$ is exact by \cite[\href{https://stacks.math.columbia.edu/tag/05R6}{Tag 05R6}]{stacks-project}.
%\end{proof}


\section[Derived Tensor Products and Derived Homs]{Derived Tensor Products and Derived $\Hom$s}

We will now work towards the definition and existence of left derived tensor products and right derived $\Hom$s. We first start off with introducing the tensor products of complexes.
\begin{definition}[Tensor products of complexes]
	\label{definition: tensor products of complexes}
	
	Let $\category{A}$ be a symmetric monoidal abelian category with all direct sums. Define a functor
	$$-\otimes_\category{A}^\bullet - \colon \homotopyCategoryOfComplexes{\category{A}} \times \homotopyCategoryOfComplexes{
		\category{A}}\to \homotopyCategoryOfComplexes{\category{A}}$$
	given on complexes $X^\bullet$ and $Y^\bullet$ in degree $n\in \Z$ by
	$$X^\bullet\otimes_\category{A}^n Y^\bullet \ldef \bigoplus_{p+q=n} X^p \otimes_\category{A} Y^q$$
	with differential
	$$d^n \ldef \bigoplus_{p+q=n} d_X^p\otimes \id_{Y^q} + (-1)^p \cdot \id_{X^p}\otimes d_Y^q.$$
	Denoting by $\mathbf{1}$ the unit object of $\category{A}$ it is clear that $-\otimes_\category{A}^\bullet -$ makes $\homotopyCategoryOfComplexes{\category{A}}$ into a symmetric monoidal category with unit given by $\mathbf{1}[0]$.
\end{definition}

We need an analogous construction for $\Hom$-sets. As there are several enrichments floating around at any given time, we give a general definition. The main specializations for us being the usual enrichment over abelian groups (as is standard for any abelian category), but also over condensed abelian groups and over condensed $A$-modules for some condensed ring $A$, recalling that both $\condensedAb$ and $\mod{A}$ are closed monoidal, \ie have an internal $\intHom$ enriching the usual $\Hom$.
\begin{definition}[Complexes of $\Hom$s]
	\label{definition: complexes of Homs}
	
	Let $\category{A}$ be an abelian category enriched over the monoidal abelian category $\category{V}$. Assume that $\category{V}$ admits all direct products. Define a functor
		$$\eHom_\category{A}^\bullet\colon \homotopyCategoryOfComplexes{\category{A}}^\op \times \homotopyCategoryOfComplexes{\category{A}}\to \homotopyCategoryOfComplexes{\category{B}}$$
	given on complexes $X^\bullet$ and $Y^\bullet$ in degree $n\in \Z$ by
		$$\eHom_\category{A}^\bullet(X^\bullet, Y^\bullet)\ldef \prod_{p+q=n} \eHom_\category{A}(X^p, Y^q)$$
	with differential
		$$d^n \ldef \prod_{p+q=n} (d_Y^q)_\ast + (-1)^{n+1} \cdot (d_X^{p-1})^\ast$$
	where as usual, $(-)_\ast$ and $(-)^\ast$ denote the images of morphisms under $\eHom_\category{A}$ when fixing the first respectively second argument.
\end{definition}

Tensor products of complexes and complexes of internal $\intHom$s inside a closed symmetric monoidal $\category{A}$ interact well and ensure that the homotopy category of complex $\homotopyCategoryOfComplexes{\category{A}}$ is closed monoidal as well.
\begin{remark}[Chain complexes are closed monoidal]
	\label{definition: chain complexes are closed monoidal}
	
	Let $\category A$ be closed symmetric monoidal abelian category with arbitrary direct sums and arbitrary direct products. Denote by $\mathbf{1}$ the unit object of $\category{A}$. Since all partial adjunctions $-\otimes_\category{A}^\bullet Y \ladj \intHom_\category{A}^\bullet(Y, -)$ for any complex $Y$ are compatible, we obtain that $\otimes_\category{A}^\bullet$ and $\intHom_\category{A}^\bullet$ make $\homotopyCategoryOfComplexes{\category{A}}$ into a closed symmetric monoidal category with unit given by $\mathbf{1}[0]$. Furthermore, we obtain an even stronger statement, namely that there is an isomorphism
	$$\intHom_\category{A}^\bullet(X \otimes_\category{A}^\bullet, Y, Z) \cong \intHom_\category{A}^\bullet(X, \intHom_\category{A}^\bullet(Y, Z))$$
	natural in all three complexes $X$, $Y$ and $Z$.
\end{remark}

We can now define what the left derived tensor product should be.
\begin{definition}[Left derived tensor products]
	\label{definition: left derived tensor products}
	
	Recall the setting of definition \ref{definition: tensor products of complexes}. Observe that for any morphism of complexes $Y\to Y'$ there is a natural transformation of functors $-\otimes_\category{A}^\bullet Y \Rightarrow -\otimes_\category{A}^\bullet Y'$ inducing a natural transformation of left derived functors $\LeftD(-\otimes_\category{A}^\bullet Y) \Rightarrow \LeftD(-\otimes_\category{A}^\bullet Y')$ in case that both exist. The \emph{left derived tensor product} is the functor
	$$-\Dotimes_\category{A} - \colon \derived{\category{A}} \times \derived{\category{A}} \to \derived{\category{A}}$$
	such that for any complex $Y$ the partial functor $-\Dotimes_\category{A} Y$ is the left derived functor $\LeftD(-\otimes_\category{A}^\bullet Y)$.
\end{definition}

In a similar fashion, we obtain the definition of right derived $\Hom$s.
\begin{definition}[Right derived $\Hom$s]
	\label{definition: right derived Homs}
	
	Recall the setting of definition \ref{definition: complexes of Homs}. Observe that for any morphism $Y\to Y'$ of complexes there is a natural transformation $\eHom_\category{A}^\bullet(-, Y) \Rightarrow \eHom_\category{A}^\bullet(-, Y')$ of functors $\homotopyCategoryOfComplexes{\category{A}}^\op \to \homotopyCategoryOfComplexes{\category{A}}$, inducing a natural transformation of right derived functors $\RightD\eHom_\category{A}^\bullet(-, Y)\Rightarrow \RightD\eHom_\category{A}^\bullet(-, Y')$ in case that both exist. The \emph{right derived $\Hom$} is the functor
		$$\RightD\eHom_\category{A}\colon \derived{\category{A}}^\op\times\derived{\category{A}}\to \derived{\category{B}}$$
	such that for any complex $Y$ the partial functor $\RightD\eHom_\category{A}(-, Y)$ is the right derived functor $\RightD\eHom_\category{A}^\bullet(-, Y)$.
\end{definition}

\begin{remark}[About the notation and choices]
	\label{remark: about the notation and choices}
	
	Recall the settings of definitions \ref{definition: left derived tensor products} and \ref{definition: right derived Homs}. One should now ask the following two questions: First, why not name the derived functors $\otimes^{\LeftD\bullet}_\category{A}$ and $\RightD\eHom_\category{A}^\bullet$ instead of dropping the $\bullet$? And secondly, why fix the second argument and not the first? After all that is an arbitrary choice.
	
	To answer the first question one has to notice that if either the first argument or second argument of $\otimes_\category{A}^\bullet$ or $\eHom_\category{A}^\bullet$ is an object of $\category{A}$ concentrated in degree $0$, let us assume the first argument $X\in\homotopyCategoryOfComplexes{\category{A}}$ is general, and the second argument is some $Y\in\category{A}$, then
	$$X\otimes_\category{A}^\bullet Y[0] \cong X\otimes_\category{A} Y[0]\text{ and }\eHom_\category{A}^\bullet(X, Y[0])\cong \eHom_\category{A}(X, Y[0])$$
	agree with the natural extensions of the functors $-\otimes_\category{A} Y$ and $\eHom_\category{A}(-, Y)$ to functors $\homotopyCategoryOfComplexes{\category{A}}\to\homotopyCategoryOfComplexes{\category{A}}$ and $\homotopyCategoryOfComplexes{\category{A}}^\op \to \homotopyCategoryOfComplexes{\category{B}}$. So in this case $- \Dotimes_\category{A} Y$ and $\RightD\eHom_\category{A}(-, Y)$ really are the derived functors expected from the notation. Because of this and since we are not really interested in $\otimes_{\category{A}}^\bullet$ and $\eHom_\category{A}^\bullet$ themselves, we simply drop the $\bullet$.
	
	To understand why we fix the second argument instead of the first, we should recall the case of modules over a discrete ring $R$. The resulting category $\mod{R}$ has both enough projectives and enough injectives. Then one might as well fix the first argument and resolve the second one by a \emph{(K-)injective} complex: Using a simple dimension shifting argument one shows by hand that the two notions agree and hence that either choice of argument yields isomorphic functors. As a consequence for either choice of argument the constructed derived functor must be also the derived functor of the other argument. An argument like this is not possible in a condensed world as there are simply not enough injectives. The author was unable to provide an alternative proof so we must make this an \assumptionPreWord.
\end{remark}

%\begin{remark}[Some enrichments]
%	Observe that in a closed symmetric monoidal abelian category $\category{A}$ with unit $\mathbf{1}$ there is an isomorphism
%		$$\eHom_\category{A}(\mathbf{1}, \intHom_\category{A}(X, Y)) \cong \eHom_\category{A}(\mathbf{1}\otimes_\category{A} X, Y) \cong \eHom_\category{A}(X, Y)$$
%	natural in the objects $X, Y\in\category{A}$, ensuring that $\intHom_\category{A}$ enriches $\eHom_\category{A}$ and hence $\Hom_\category{A}$. Thus the natural extension of $\eHom_\category{A} \colon \category{A}\to \Ab$ to a functor $\complexes{\category{A}}\to\complexes{\Ab}$ provides an isomorphism
%		$$\eHom_\category{A}(\mathbf{1}, -) \circ \intHom_\category{A}^\bullet \cong \eHom_\category{A}^\bullet.$$
%\end{remark}

Using out previous results on $K$-projectivity we can now ensure the existence of the left derived tensor product.
\begin{corollary}[Existence of the left derived tensor product]
	Assume further in the setting of definition \ref{definition: left derived tensor products} that $\category{A}$ has enough projectives and that filtered colimits in $\category{A}$ are exact. Then the left derived tensor product $\Dotimes_\category{A}$ exists.
\end{corollary}
\begin{proof}
	This is direct consequence of lemma \ref{lemma: existence of left derived functors} once one knows that there are enough $K$-projective resolutions due to proposition \ref{proposition: enough projectives on steroids}.
\end{proof}

Similarly, we obtain the existence of right derived $\Hom$s.
\begin{corollary}[Existence of right derived $\Hom$s]
	\label{corollary: existence of right derived Homs}
	
	Assume further in the setting of definition \ref{definition: right derived Homs} that $\category{A}$ has enough projectives and that filtered colimits in $\category{A}$ are exact. Then the right derived functor $\RightD\eHom_\category{A}$ exists.
\end{corollary}
\begin{proof}
	Observe that for any complex $Y$ the right derived functor of $\eHom_\category{A}^\bullet(-, Y)\colon \homotopyCategoryOfComplexes{\category{A}}^\op \to \homotopyCategoryOfComplexes{\category{B}}$ is equivalently the left derived functor of $\eHom_\category{A}^\bullet(-, Y)\colon \homotopyCategoryOfComplexes{\category{A}} \to \homotopyCategoryOfComplexes{\category{B}}^\op\equiv \homotopyCategoryOfComplexes{\category{B}^\op}$. Then once again one knows that there are enough $K$-projective resolutions due to proposition \ref{proposition: enough projectives on steroids} giving the existence of the left derived functor by \ref{lemma: existence of left derived functors}.
\end{proof}


Denoting by $\eHom_\category{A}$ then enrichment of $\Hom_\category{A}$ to abelian groups, we see that $\eRHom_\category{A}$ actually enriches $\eHom_\derived{\category{A}}$. In particular in the derived category, to each $\Hom$-set there is a corresponding complex of abelian groups.
\begin{lemma}[$\eRHom$ enriches $\Hom$-sets in the derived category]
	\label{lemma: eRHom enriches Hom-sets in the derived category}
	
	Let $\category{A}$ be an abelian category that admits all direct products, has enough projectives and such that filtered colimits are exact. Then $\eRHom_\category{A}\colon \derived{\category{A}}^\op \times\derived{\category{A}}\to\derived{\Ab}$ exists by corollary \ref{corollary: existence of right derived Homs}, and there is a natural isomorphism
	$$H^0\eRHom_\category{A} \cong \eHom_\derived{\category{A}},$$
	providing a natural enrichment of $\eHom_\derived{\category{A}}$ by $\eRHom_\category{A}$.
\end{lemma}
\begin{proof}
	Let $X^\bullet$ and $Y^\bullet$ be complexes. Choosing a left resolution of $X^\bullet$ by a $K$-projective complex $P^\bullet$ we see that
	$$H^0\eRHom_\category{A}(X^\bullet, Y^\bullet) \cong H^0 \eHom_\category{A}^\bullet(P^\bullet, Y^\bullet) \cong \eHom_\homotopyCategoryOfComplexes{\category{A}}(P^\bullet, Y^\bullet).$$
	As $P^\bullet$ is $K$-projective we obtain by lemma \ref{lemma: morphism out of a K-projective complex} that
	$$\eHom_\homotopyCategoryOfComplexes{\category{A}}(P^\bullet, Y^\bullet)\cong \eHom_\derived{\category{A}}(P^\bullet, Y^\bullet) \cong \eHom_\derived{\category{A}}(X^\bullet, Y^\bullet),$$
	proving the claim.
\end{proof}

\section{Deriving Adjunctions}

In this section we will study when adjunctions are stable under derivation, in particular enriched adjunctions. If there is no enrichment we have the following general result.
\begin{lemma}[Adjunctions are stable under derivation, {\cite[\href{https://stacks.math.columbia.edu/tag/09T5}{Tag 09T5}]{stacks-project}}]
	\label{lemma: adjunctions are stable under derivation}
	
	Let $\category{A}$ and $\category{B}$ be abelian categories and $F\colon \category{A} \to \category{B}$ be left adjoint to $G\colon \category{B}\to\category{A}$. If $\LeftD F$ and $\RightD G$ exist, then there are isomorphisms
	$$\Hom_\derived{\category{B}}(\LeftD F(X^\bullet), Y^\bullet) \cong \Hom_\derived{\category{A}}(X^\bullet, \RightD G(Y^\bullet))$$
	natural in the complexes $X^\bullet \in \derived{\category{A}}$ and $Y^\bullet \in \derived{\category{B}}$.
\end{lemma}

An immediate consequence is that a closed monoidal structure on a (well-behaved) abelian category translates to the derived setting.
\begin{lemma}[Deriving a closed monoidal structure]
	\label{lemma: deriving a closed monoidal structure}
	
	Let $\category{A}$ be a closed symmetric monoidal abelian category. Assume that $\category{A}$ is \grothendieckaxiom{$5$} and has enough projectives. Then there is an isomorphism
		$$\Hom_\derived{\category{A}}(X\Dotimes_\category{A}Y, Z)\cong \Hom_\derived{\category{A}}(X, \intRHom_\category{A}(Y, Z))$$
	natural in the complexes $X$, $Y$ and $Z$ making $\derived{\category{A}}$ into a closed symmetric monoidal category. Furthermore, there even is a natural isomorphism
		$$\intRHom_\category{A}(X\Dotimes_\category{A}Y, Z)\cong \intRHom_\category{A}(X, \intRHom_\category{A}(Y, Z))$$
	of enrichments of $\Hom_\derived{\category{A}}$.
\end{lemma}
\begin{proof}
	By lemma \ref{lemma: adjunctions are stable under derivation} we obtain that for any complex $X$ the two functors $- \Dotimes_\category{A} X$ and $\intRHom_\category{A}(X, -)$ are adjoint. As these adjunctions for various $X$ are compatible it is easily verified that $\derived{\category{A}}$ is closed symmetric monoidal with unit $\mathbf{1}[0]$, the unit of $\category{A}$ concentrated in degree $0$.
	
	It is now purely formal that the adjunctions $-\Dotimes_\category{A}X\ladj \intRHom_\category{A}(X, -)$ enrich. Indeed, in any closed monoidal category $\category{V}$ and $X, Y, Z\in \category{V}$ we find that
	\begin{equation}
		\scalebox{0.95}{\parbox{\textwidth}{
			\begin{align*}
				\Hom(-, \intHom(X\otimes Y, Z))&\cong \Hom(-\otimes X\otimes Y, Z)\\
				&\cong \Hom(-\otimes X, \intHom(Y, Z))\cong \Hom(-, \intHom(X, \intHom(Y, Z)))
			\end{align*}
		}}
	\end{equation}
	implying by Yoneda's lemma that there is a natural isomorphism $\intHom(X\otimes Y, Z)\cong \intHom(X,\intHom(Y, Z))$ implying that $-\otimes X\ladj_\category{V} \intHom(X, -)$ for all $X \in \category{V}$.
\end{proof}

More generally we are also interested in translating enriched adjunctions $F\ladj_\category{V} G$ to the derived category. In this thesis the main application of this is the case of scalar extension and restriction where the adjunction is enriched over condensed abelian groups (or over the respective module categories). This however is somewhat precarious and the author only managed to prove that an enriched adjunction translates under quite strong assumptions (tailored to this specific use case). In the condensed setting we find that scalar extension maps compact projective generators to compact projective generators all while scalar restriction is exact. This helps a lot in applying Spaltenstein's techniques.
\begin{lemma}[Enriched adjunctions are stable under derivation]
	\label{lemma: enriched adjunctions are stable under derivation}
	
	Let $\category{A}$ and $\category{B}$ be abelian categories, both enriched over the monoidal abelian category $\category{V}$. Assume that all three categories are \grothendieckaxiom{$5$} with enough projectives. Let $F\colon \category{A} \to \category{B}$ and $G\colon \category{B}\to\category{A}$ be a pair of additive functors such that $F\ladj_\category{V} G$ is a $\category{V}$-enriched adjunction, such that $F$ maps projectives to projectives and such that $G$ is exact. Then in $\derived{\category{V}}$ there are isomorphisms
	$$\eRHom_\category{B}(\LeftD F(X^\bullet), Y^\bullet) \cong \eRHom_\category{A}(X^\bullet, \RightD G(Y^\bullet))$$
	natural in the complexes $X^\bullet \in \derived{\category{A}}$ and $Y^\bullet \in \derived{\category{B}}$.
\end{lemma}
\begin{proof}
	From $\eHom_\category{B}(F(-), -)\cong \eHom_\category{A}(-, G(-))$ it follows formally that
	$$\eHom_\category{B}^\bullet(\tilde F(-), -)\cong \eHom^\bullet_\category{A}(-, \tilde G(-))$$
	where $\tilde F$ and $\tilde G$ are the natural extensions of $F$ and $G$ to accommodate complexes. Denote by $\mathfrak{P}_\category{A}$ the class of bounded to the right complexes in $\category{A}$ which are term-wise projective and denote similarly $\mathfrak{P}_\category{B}$. Recall from \cite[Def. 2.6]{spaltenstein} the notion of a \emph{$\mathfrak{P}_\category{A}$-special direct system}. Suppose that $(P_n^\bullet)_{n\in E}$ is such a $\mathfrak{P}_\category{A}$-special direct system with colimit $P^\bullet$. Then $\tilde F(P^\bullet_n) \in \mathfrak{P}_\category{B}$ for any $n\in E$ as $F$ maps projectives to projectives by assumption. Now as $F$ is additive, it preserves split exact sequences in $\category{A}$ and hence $\tilde F$ preserves \emph{semi-split sequences} (see \cite[Def. 0.5]{spaltenstein}) of complexes in particular $\tilde F(P^\bullet_{n-1})\to\tilde F(P^\bullet_n)$ is injective as well, for any $n\in E$ which has a predecessor $n-1$. We thus obtain that $(\tilde F(P^\bullet_n))_{n\in E}$ is a $\mathfrak{P}_\category{B}$-special direct system. As $\tilde F$ is left adjoint to $\tilde G$ and thus preserves colimits, we obtain that $\tilde F(P^\bullet) = \tilde F(\colimit_n P^\bullet_n) \cong \colimit_n \tilde F(P^\bullet_n)$ is a colimit of a $\mathfrak{P}_\category{B}$-special direct system and hence must be $K$-projective.
	
	Let $X^\bullet \in \derived{\category{A}}$ and $Y^\bullet \in \derived{\category{B}}$. Choose a left resolution $P^\bullet \to X^\bullet$ by a $K$-projective complex $P^\bullet$ that is the colimit of a $\mathfrak{P}_\category{A}$-special direct system. Then by the above argument $\tilde F(P^\bullet)$ is $K$-projective and we obtain that
	\begin{align*}
		\eRHom_\category{B}(\LeftD F(X^\bullet), Y^\bullet)
		&\cong\eRHom_\category{B}(\tilde F(P^\bullet), Y^\bullet)\\
		&\cong\eHom_\category{B}^\bullet(\tilde F(P^\bullet), Y^\bullet)\\
		&\cong \eHom_\category{A}^\bullet(P^\bullet, \tilde G(Y^\bullet)) = \eRHom_\category{A}(X^\bullet, \RightD G(Y^\bullet))
	\end{align*}
	proving the claim.
\end{proof}

%The case of a general closed symmetric monoidal category (as for example condensed abelian groups or condensed modules over some condensed ring) still eludes the author. In case that the projectives are also flat (which is the case for such categories of condensed objects) and are stable under tensor products (which is not(!) the case for these categories) we obtain the following result.
%\begin{lemma}
%	Suppose that $\category{A}$ is closed symmetric monoidal, admits all direct sums and direct products, has enough projectives and that filtered colimits in $\category{A}$ are exact. Suppose further that $\category{A}$ has enough projectives that are additionally flat and that tensor products of projectives are projective. Then there is an isomorphism
%	$$\intRHom_\category{A}(X\Dotimes_\category{A} Y, Z) \cong \intRHom_\category{A}(X, \intRHom_\category{A}(Y, Z))$$
%	natural in the complexes $X$, $Y$ and $Z$. This makes $\derived{\category{A}}$ into a closed symmetric monoidal category.
%\end{lemma}
%\begin{proofsketch}
%	\begin{todo}
%		Content
%		Need to argue that special direct limits of $K$-flats are again $K$-flat. This is \cite[Prop. 2.8]{spaltenstein}
%	\end{todo}
%	Choose $K$-projective left resolutions $P_X$ of $X$ resp. $P_Y$ of $Y$ that both are sequential colimits of systems of complexes $(P_{X, n})_{n\ge 1}$ resp. $(P_{Y,n})_{n\ge -1}$ where for each $n\ge -1$ the complexes $P_{X, n}$ resp. $P_{Y, n}$ are term-wise projective and flat. Now by definition we have $X\Dotimes_\category{A} Y = P_X\otimes_\category{A}^\bullet Y$. Now since all $P_{X,n}$ consist of flat objects, each such complex and hence their colimit are \emph{$K$-flat} that is for each acyclic complex $A$ the complex $P_X\otimes_\category{A}^\bullet A$ is acyclic as well. This implies that $P_X\otimes_\category{A}^\bullet -$ preserves quasi-isomorphisms. Indeed if $s\colon Y\to Y'$ is an quasi-isomorphism then its cone $C$ is acyclic. We then have a distinguished triangle
%	$$P_X\otimes_\category{A}^\bullet Y \to P_X\otimes_\category{A}^\bullet Y'\to P_X\otimes_\category{A}^\bullet C \to P_X\otimes_\category{A}^\bullet Y[1]$$
%	where $P_X\otimes_\category{A}^\bullet C$ is acyclic by $K$-flatness of $P_X$. Thus $P_X\otimes_\category{A}^\bullet Y\to P_X\otimes_\category{A}^\bullet Y'$ is a quasi-isomorphism as well. In particular this applies to the quasi-isomorphism $P_Y\to Y$. We thus obtain that $P_X\otimes_\category{A}^\bullet Y$ is quasi-isomorphic to $P_X\otimes_\category{A}^\bullet P_Y$. Now observe that
%	$$P_X\otimes_\category{A}^\bullet P_Y \cong \colimit_{n, m\ge -1} P_{X,n}\otimes_\category{A}P_{Y, n}.$$
%	But each $P_{X,n}\otimes_\category{A}^\bullet P_{Y,n}$ is a complex of projectives, as tensor products of projectives are again projective, yielding that their filtered colimit $P_X\otimes_\category{A}^\bullet P_Y$ is $K$-projective as it can be written as a special direct system of the diagonal $m=n$. Hence
%	\begin{align*}
%		\intRHom_\category{A}(X\Dotimes_\category{A} Y, Z)
%		&\cong \intRHom_\category{A}(P_X\otimes_\category{A}^\bullet P_Y, Z)\\
%		&\cong \intHom_\category{A}^\bullet(P_X \otimes_\category{A}^\bullet P_Y, Z)\\
%		&\cong \intHom_\category{A}^\bullet(P_X, \intHom_\category{A}^\bullet(P_Y, Z)) \cong \intRHom_\category{A}(X, \intRHom_\category{A}(Y, Z))
%	\end{align*}
%	which proves the claim.
%\end{proofsketch}

%In the general case we now have to make the following assumption.
%\begin{assumption}[Deriving a closed symmetric monoidal category]
%	\label{assumption: deriving a closed symmetric monoidal category}
%	
%	Let $\category{A}$ be a symmetric monoidal abelian category such that $\Dotimes_\category{A}$ and $\intRHom_\category{A}$ exist. Then there are isomorphisms
%		$$\intRHom_\category{A}(X\Dotimes_\category{A} Y, Z)\cong \intRHom_\category{A}(X, \intRHom_\category{A}(Y, Z))$$
%	natural in the complexes $X$, $Y$ and $Z$. This makes $\derived{\category{A}}$ into a closed symmetric monoidal category.
%\end{assumption}
%
%%\begin{assumption}
%%	Let $A$ be a condensed ring. Write $\derived{A}$ for the derived category of $\mod A$. Both the left derived tensor product $-\Dotimes_A -$ and the right derived $\intRHom_A$ exist and there is an isomorphism
%%	$$\intRHom_A(L^\bullet\Dotimes_A M^\bullet, N^\bullet)\cong \intRHom_A(L^\bullet, \intRHom_A(M^\bullet, N^\bullet))$$
%%	natural in the complexes $L^\bullet$, $M^\bullet$ and $N^\bullet$.
%%\end{assumption}

While the previous results cover many cases of interest, occasionally we need the following more general assumption.
\begin{assumption}[Derived adjunctions are enriched]
	\label{assumption: derived adjunctions are enriched}
	
	Let $\category{A}$ and $\category{B}$ be abelian categories enriched over the monoidal abelian category $\category{V}$. Assume that for each of the three categories the associated derived category is locally small. If $F\colon \derived{\category{A}} \to \derived{\category{B}}$ if left adjoint to $G\colon \derived{\category{B}}\to\derived{\category{A}}$ then in $\derived{\category{V}}$ there are isomorphisms
		$$\eRHom_\category{B}(F(X^\bullet), Y^\bullet) \cong \eRHom_\category{A}(X^\bullet, G(Y^\bullet))$$
	natural in the complexes $X^\bullet \in \derived{\category{A}}$ and $Y^\bullet \in \derived{\category{B}}$ providing a natural enrichment of the adjunction.
\end{assumption}

\newpage
\section{(Commuting with) Homotopy (Co-)Limits}
 
It is often very convenient to break down an argument about an unbounded complex to a left or right bounded one. Recall the several notions of truncation of complexes from \cite[\href{https://stacks.math.columbia.edu/tag/0118}{Tag 0118}]{stacks-project}.
\begin{lemma}[A complex as a (co)limit of its truncations]
	\label{lemma: a complex as a (co)limit of its truncations}
	
	Let $\category{A}$ be an abelian category and $C^\bullet \in \complexes{\category{A}}$ any complex.
	\begin{itemize}
		\item
		The stupid truncations $\sTruncation{\le n}C^\bullet$ with the natural projections $\sTruncation{\le m} \to \sTruncation{\le n}$ for $n\le m$ form a sequential inverse system. If $\category{A}$ has sequential limits then $C^\bullet \to \limit_{n\in \N} \sTruncation{\le n} C^\bullet$ is an isomorphism.
		
		\item
		The stupid truncations $\sTruncation{\le n}C^\bullet$ with the natural inclusions $\sTruncation{\le n} \to \sTruncation{\le m}$ for $n\le m$ form a sequential direct system. If $\category{A}$ has sequential colimits then $\colimit_{n\in \N} \sTruncation{\le n}C^\bullet \to C^\bullet$ is an isomorphism.
		
		\item
		The canonical truncations $\cTruncation{\le n}C^\bullet$ with the natural inclusions $\cTruncation{\le n} \to \cTruncation{\le m}$ for $n\le m$ form a sequential direct system. If $\category{A}$ has sequential colimits then $\colimit_{n\in \N} \cTruncation{\le n}C^\bullet \to C^\bullet$ is an isomorphism.
		
		\item
		The stupid truncations $\sTruncation{\ge -n}C^\bullet$ with the natural inclusions $\sTruncation{\ge -n}C^\bullet \to \sTruncation{\ge -m}C^\bullet$ for $n\le m$ for a sequential direct system. If $\category{A}$ has sequential colimits them $\colimit_{n\in \N} \sTruncation{\ge -n} C^\bullet \to C^\bullet$ is an isomorphism.
	\end{itemize}
\end{lemma}
\begin{shortproof}
	Clear since in each fixed degree $n\in \Z$ each diagram is eventually constant.
\end{shortproof}

Recall the definition of (sequential) homotopy limits from \cite[\href{https://stacks.math.columbia.edu/tag/08TC}{Tag 08TC}]{stacks-project} and of (sequential) homotopy colimits from \cite[\href{https://stacks.math.columbia.edu/tag/090Z}{Tag 090Z}]{stacks-project}.
\begin{lemma}[Ordinary sequential colimits are homotopy colimits, {\cite[\href{https://stacks.math.columbia.edu/tag/0949}{Tag 0949}]{stacks-project}}]
	\label{lemma: oridnary sequential colimits are homotopy limits}
	
	Let $\category{A}$ be an abelian category such that sequential colimits exist and are exact (for example if $\category{A}$ is  \grothendieckaxiom{$5$}). For any sequential diagram $n\mapsto X_n^\bullet$ of complexes in $\category{A}$ the colimit of complexes $\colimit_{n\in \N} X_n^\bullet$ is a homotopy colimit of $n\mapsto X_n^\bullet$ in $\derived{\category{A}}$.
\end{lemma}
This implies that any of the colimits in lemma \ref{lemma: a complex as a (co)limit of its truncations} is actually a homotopy colimit. For non-sequential diagrams the definition of homotopy (co)limits is more inexplicit.
\begin{definition}[Homotopy limits and colimits]
	Let $\category{A}$ be an abelian category and $\category{I}$ be a small category. Consider the abelian category $\functorCategory{\category{I}}{\category{A}}$ of $\category{I}$-shaped diagrams in $\category{A}$. Observe that we have additive functors
		$$\limit\colon \functorCategory{\category{I}}{\category{A}} \to \category{A}\text{ and }\colimit\colon \functorCategory{\category{I}}{\category{A}}\to\category{A}$$
	if $\category{A}$ is complete respectively cocomplete. Let $D\colon \category{I} \to \complexes{\category{A}}$ be a diagram. In case that the right derived functor $\RightD\limit$ is defined at $D$ we call its image
		$$\homotopyLimit D \ldef (\RightD\limit)(D) \in \derived{\category{A}}$$
	the \emph{homotopy limit of $D$}. Similarly, if the left derived functor $\LeftD\colimit$ is defined at $D$ we call its image
		$$\homotopyColimit D\ldef (\LeftD \colimit)(D) \in \derived{\category{A}}$$
	the \emph{homotopy colimit of $D$}.
\end{definition}

%\begin{assumption}
%	\begin{todo}
%		Do we need this?
%	\end{todo}
%	Suppose that $\category{A}$ is an abelian category enriched over the monoidal abelian category $\category{V}$ such that both category are \grothendieckaxiom{$5$} and have enough projectives. If $n\mapsto Y_n$ is any sequential diagram of complexes then the morphism
%	$$\eRHom_A(X, \homotopyLimit_{n\in \N} Y_n) \to \homotopyLimit_{n\in \N} \eRHom_A(X, Y_n)$$
%	in $\derived{\category{A}}$ is an isomorphism.
%\end{assumption}

\begin{remark}
	If the abelian category $\category{A}$ is \grothendieckaxiom{$5$} and $\category{I}$ is filtered then by definition $\colimit\colon \functorCategory{\category{I}}{\category{A}} \to \category{A}$ is exact and hence $\homotopyColimit D \cong \colimit D$ for any diagram $D\colon \category{I}\to\complexes{\category{A}}$.
\end{remark}

We now find that derived $\Hom$s maps (sequential) homotopy colimits to homotopy limits in the first argument, as expected.
\begin{lemma}[Right derived Hom \& homotopy colimits in the first argument]
	\label{lemma: right derived Hom and homotopy colimits in the first argument}
	
	Let $\category{A}$ be an abelian category enriched over the monoidal abelian category $\category{V}$. Assume that both categories are \grothendieckaxiom{$5$} and have enough projectives. Let $Y$ be a complex and assume that $\eRHom_\category{A}(-, Y)$ maps direct sums to direct products. For any sequential diagram $n\mapsto X_n$ of complexes the complex $\eRHom_\category{A}(\colimit_{n\in \N} X_n, Y)$ is a homotopy limit of $n\mapsto \eRHom_\category{A}(X_n, Y)$ in $\derived{\category{V}}$. We thus may write
		$$\eRHom_\category{A}(\colimit_{n\in \N} X_n, Y) \cong \homotopyLimit_{n\in \N} \eRHom_\category{A}(X_n ,Y).$$
\end{lemma}
\begin{proof}
	Since $\colimit_{n\in \N} X_n$ is a homotopy colimit of $n\mapsto X_n$ there is a distinguished triangle
		$$\bigoplus_{n\in \N} X_n\to \bigoplus_{n\in \N} X_n\to \colimit_{n\in \N} X_n \to \bigoplus_{n\in \N} X_n[1].$$
	Application of $\eRHom_\category{A}(-, Y)$ yields that
		$$\eRHom_\category{A}(\colimit_{n\in \N} X_n, Y) \to \eRHom_\category{A}(\bigoplus_{n\in \N} X_n, Y)\to \eRHom_\category{A}(\bigoplus_{n\in \N} X_n, Y)\to (...)[1]$$
	is a distinguished triangle as well. As $\eRHom_\category{A}(-, Y)$ maps direct sums to products we obtain a morphism of distinguished triangles
	$$\begin{tikzcd}[column sep=1em, scale cd=0.9]
		\eRHom_\category{A}(\colimit_{n\in \N} X_n, Y) \ar[r]\ar[d, equal] &\eRHom_\category{A}(\bigoplus_{n\in \N} X_n, Y) \ar[r]\ar[d] &\eRHom_\category{A}(\bigoplus_{n\in \N} X_n, Y) \ar[r]\ar[d] &(...)[1]\\
		\eRHom_\category{A}(\colimit_{n\in \N} X_n, Y) \ar[r] & \prod_{n\in \N} \eRHom_\category{A}(X_n, Y) \ar[r] &\prod_{n\in \N} \eRHom_\category{A}(X_n, Y) \ar[r] &(...)[1].
	\end{tikzcd}$$
	Hence $\eRHom_\category{A}(\colimit_{n\in \N} X_n ,Y)$ is a homotopy limit of $n\mapsto \eRHom_\category{A}(X_n, Y)$.
\end{proof}

Completely analogous is of course the commutation with homotopy limits in the second argument.
\begin{lemma}[Right derived Hom \& homotopy limits in the second argument]
	\label{lemma: right derived Hom and homotopy limits in the second argument}
	
	Let $\category{A}$ be an abelian category enriched over the monoidal abelian category $\category{V}$. Assume that both categories are \grothendieckaxiom{$5$} and have enough projectives. Let $X$ be a complex and assume that $\eRHom_\category{A}(X, -)$ commutes with direct products. For any sequential diagram $n\mapsto Y_n$ of complexes we find that
		$$\eRHom_\category{A}(X, \homotopyLimit_{n\in \N} Y_n) \cong \homotopyLimit_{n\in \N} \eRHom_\category{A}(X, Y_n)$$
	in $\derived{\category{V}}$.
\end{lemma}
\begin{proof}
	Since $\homotopyLimit_{n\in \N} Y_n$ is a homotopy limit of $n\mapsto Y_n$ there is a distinguished triangle
		$$\homotopyLimit_{n\in \N} Y_n\to \prod_{n\in \N} Y_n\to  \prod_{n\in \N} Y_n \to (\cdots)[1].$$
	Application of $\eRHom_\category{A}(X, -)$ yields that
		$$\eRHom_\category{A}(X, \homotopyLimit_{n\in \N} Y_n) \to \eRHom_\category{A}(X, \prod_{n\in\N} Y_n)\to \eRHom_\category{A}(X, \prod_{n\in \N} Y_n)\to (...)[1]$$
	is a distinguished triangle as well. As $\eRHom_\category{A}(X, -)$ maps products to products we obtain a morphism of distinguished triangles
	$$\begin{tikzcd}[column sep=1em, scale cd=0.9]
		\eRHom_\category{A}(X, \homotopyLimit_{n\in \N} Y_n) \ar[r]\ar[d, equal] &\eRHom_\category{A}(X, \prod_{n\in \N} Y_n) \ar[r]\ar[d] &\eRHom_\category{A}(X, \prod_{n\in \N} Y_n) \ar[r]\ar[d] &(...)[1]\\
		\eRHom_\category{A}(X, \homotopyLimit_{n\in \N} Y_n) \ar[r] & \prod_{n\in \N} \eRHom_\category{A}(X, Y_n) \ar[r] &\prod_{n\in \N} \eRHom_\category{A}(X, Y_n) \ar[r] &(...)[1]
	\end{tikzcd}$$
	implying that $\eRHom_\category{A}(X, \homotopyLimit_{n\in \N} Y_n)$ is a homotopy limit of $n\mapsto \eRHom_\category{A}(X, Y_n)$.
\end{proof}

\begin{remark}[Compatibility of enriched $\Hom$s with direct sums and products]
	\label{remark: compatibility of enriched Homs with direct sums and products}
	
	That $\eRHom_\category{A}$ maps direct sums to direct products is often true. Assume for example that the category $\category{D}$ is closed symmetric monoidal (as is the case for a closed symmetric monoidal abelian category $\category{A}$ where $\Dotimes_\category{A}$ and $\intRHom_\category{A}$ exist and are (partially) adjoint, see \ref{lemma: deriving a closed monoidal structure}). Then for $X_i$ and $Y$ in $\category{D}$ we obtain that
	\begin{align*}
		\Hom(-,\intHom(\bigoplus_i X_i, Y)) &\cong \Hom(-\otimes \bigoplus_i X_i, Y)\\
		&\cong \Hom(\bigoplus_i X_i, \intHom(-, Y))\\
		&\cong \prod_i \Hom(X_i, \intHom(-, Y))\\
		&\cong \prod_i \Hom(-, \intHom(X_i, Y)) \cong \Hom(-, \prod_i \intHom(X_i, Y))
	\end{align*}
	as $\Hom$ certainly maps colimits to limits in the first argument. By Yoneda we thus find that $\intHom(\bigoplus_i X_i, Y) \cong \prod_i \intHom(X_i, Y)$ as desired -- although technically one should check that this is the correct comparison map. Similarly, one shows that $\intHom(X, \prod_i Y_i) \cong \prod_i \intHom(X, Y_i)$ as well.
\end{remark}

The commutation with general homotopy colimits is not so clear and for now we assume it given.
\begin{assumption}[Right derived $\Hom$ maps colimits to limits]
	\label{assumption: right derived Hom maps colimits to limits}
	
	Let $\category{A}$ be an abelian category enriched over the monoidal abelian category $\category{V}$ such that both categories are \grothendieckaxiom{$5$} with enough projectives. Suppose that $i\mapsto X_i^\bullet$ is a filtered system of complexes, then the natural morphism
	$$\eRHom_\category{A}(\colimit_i X_i^\bullet, Y^\bullet) \to \homotopyLimit_i \eRHom_\category{A}(X_i^\bullet, Y^\bullet)$$
	in $\derived{\category{V}}$ is an isomorphism.
\end{assumption}
%\begin{proof}
%	$\eRHom$ maps homotopy colimits to homotopy limits? Grothendieck iso? Composition of Hom with colimit gives composition of derived Hom with derived colimit. Since Hom commutes with colimits then the other way around works too
%\end{proof}

\section{Derived Compactness}

We will need slight reformulations of compactness and generation when dealing with triangulated categories. Recall that a triangulated category that has all direct sums already has all (homotopy) colimits, this can be found in \cite[\href{https://stacks.math.columbia.edu/tag/0A5K}{Tag 0A5K}]{stacks-project}. The reformulation of compactness for triangulated categories will only talk about direct sums. 
\begin{definition}[(Derived) compactness, {\cite[\href{https://stacks.math.columbia.edu/tag/07LS}{Tag 07LS}]{stacks-project}}]
	
	Let $\category{D}$ be a triangulated category with arbitrary direct sums. An object $K\in \category{D}$ is called \emph{(derived) compact} if for all direct sums $\bigoplus_{i \in I} X_i$ in $\category{D}$ the natural morphism
	$$\bigoplus\nolimits_{i\in I} \Hom_\category{D}(K, X_i) \to \Hom_\category{D}(K, \bigoplus\nolimits_{i\in I} X_i)$$
	is an isomorphism.
\end{definition}

\begin{example}[Bounded complexes of compact projectives are derived compact]
	\label{example: bounded complexes of compact projectives are derived compact}
	
	Let $\category{A}$ be an \grothendieckaxiom{$5$} abelian category generated by compact projectives. Observe that for any compact projective $P$ of $\category{A}$ the complex $P[0]$ is $K$-projective so that for every complex $X$ we find by lemma \ref{lemma: morphism out of a K-projective complex} that
		$$\Hom_\homotopyCategoryOfComplexes{\category{A}}(P[0], X)\cong \Hom_\derived{\category{A}}(P, X).$$
	But as $P$ is compact this directly implies that $\Hom_\derived{\category{A}}(P[0], -)$ commutes with arbitrary direct sums as we can verify this in the homotopy category of complexes. Hence, for each compact projective $P$ the complex $P[0]$ is derived compact. Furthermore, using an induction on the length (the difference between the largest degree after which there is no entry and the smallest degree before which there is no entry), one easily checks that this implies that any \emph{bounded} complex $P^\bullet$ of which each term is compact projective is derived projective as well.
\end{example}

While in general, contrary to notion of compactness introduced in definition \ref{definition: regular projectivity and compactness}, we might not commute with all filtered (homotopy) colimits (at least it is not known to the author if it holds true), we still can be sure that we commute with sequential colimits.
\begin{lemma}[Commutation with sequential colimits]
	\label{lemma: commutation with sequential colimits}
	
	Let $\category{A}$ be an abelian category in which sequential colimits exist and are exact (for example if $\category{A}$ is \grothendieckaxiom{5}). Suppose that $\N \to \complexes{\category{A}}, n\mapsto X_n^\bullet$ is a sequential diagram in $\complexes{\category{A}}$ and that $K^\bullet \in \derived{\category{A}}$ is derived compact. Then the natural morphism
	$$\colimit_{n\in \N} \Hom_\derived{A}(K^\bullet, X_n^\bullet) \to \Hom_\derived{A}(K^\bullet, \colimit_{n\in \N} X_n^\bullet)$$
	is an isomorphism. Even more, the natural morphism
	$$\colimit_{n\in \N} \eRHom_\category{A}(K^\bullet, X_n^\bullet) \to \eRHom_\category{A}(K^\bullet, \colimit_{n\in \N} X_n^\bullet)$$
	in $\derived{\Ab}$ is an isomorphism, providing a natural enrichment.
\end{lemma}
\begin{proof}
	The first isomorphism is precisely \cite[\href{https://stacks.math.columbia.edu/tag/094A}{Tag 094A}]{stacks-project} after noticing that (sequential) homotopy colimits and ordinary colimits agree by lemma \ref{lemma: oridnary sequential colimits are homotopy limits}. For the second isomorphism, we can check this in each degree $m\in \Z$. Indeed, as 
	$$H^m \eRHom_\category{A}(-, -)\cong H^0\eRHom_\category{A}(-, -[m]) \cong \eHom_\derived{\category{A}}(-, -[m])$$
	this follows from the first case, as then in degree $m$ the induced morphism in cohomology
	$$\colimit_{n\in \N} \eHom_\derived{A}(K^\bullet, X_n^\bullet[m])\to \eHom_\derived{A}(K^\bullet, \colimit_{n\in \N} X_n^\bullet[m])$$
	is an isomorphism since $\colimit_{n\in \N}$ is exact and hence commutes with $H^n$.
\end{proof}

%One can extend the above lemma to iterated sequential colimits.
%\begin{remark}[Commutation with iterated sequential colimits]
%	\label{remark: commutation with iterated sequential colimits}
%	
%	In the setting of lemma \ref{lemma: commutation with sequential colimits}, by applying said lemma to the category of functors $\functorCategory{\N}{\category{A}}$ (which still is abelian, which still has all sequential colimits and in which sequential colimits are still exact) with its exact inclusion $\category{A} \to \functorCategory{\N}{\category{A}}, X\mapsto (n\mapsto X)$, we see that we can iterate on sequential colimits. We obtain that for any diagram $\N\times \N\to \complexes{\category{A}}, (m,n)\mapsto X_{m,n}^\bullet$ the natural morphisms
%	$$\colimit_{m,n\in \N} \Hom_\derived{\category{A}}(K^\bullet, X_{m,n}^\bullet)\to \Hom_\derived{\category{A}}(K^\bullet, \colimit_{m,n\in \N} X_{m, n}^\bullet)$$
%	and
%	$$\colimit_{m,n\in \N} \eRHom_\category{A}(K^\bullet, X_{m,n}^\bullet)\to \eRHom_\category{A}(K^\bullet, \colimit_{m,n\in \N} X_{m, n}^\bullet)$$
%	are isomorphisms as well. 
%\end{remark}


In case that the triangulated category in question is the derived category of a monoidally enriched abelian category we can generalize the notion of derived compactness.
\begin{definition}[Enriched compactness]
	\label{definition: enriched compactness}
	
	Let $\category{A}$ be an abelian category enriched over the monoidal abelian category $\category{V}$. Assume that both categories are \grothendieckaxiom{$5$} and have enough projectives. An object $K \in \derived{\category{A}}$ is called \emph{($\category{V}$-)enriched (derived) compact} if for all direct sums $\bigoplus_{i\in I} X_i$ in $\derived{\category{A}}$ the natural morphism
		$$\bigoplus_{i\in I} \eRHom_\category{A}(K, X_i) \to \eRHom_\category{A}(K, \bigoplus_{i\in I} X_i)$$
	in $\derived{\category{V}}$ is an isomorphism.
\end{definition}


For enriched compacts we can now show commutation with sequential homotopy colimits.
\begin{lemma}[Enriched compact objects preserve sequential homotopy colimits]
	\label{lemma: enriched compact objects preserve sequential homotopy colimits}
	
	Let $\category{A}$ be an abelian category enriched over the monoidal abelian category $\category{V}$. Assume that both categories are \grothendieckaxiom{$5$} and have enough projectives. If $K$ is enriched compact and $n\mapsto X_n^\bullet$ is a sequential diagram of complexes then the natural morphism
		$$\colimit_{n\in \N} \eRHom_\category{A}(K^\bullet, X_n^\bullet) \to \eRHom_\category{A}(K^\bullet, \colimit_{n\in \N} X_n^\bullet)$$
	in $\derived{\category{V}}$ is an isomorphism.
\end{lemma}
\begin{proof}
	Consider the defining distinguished triangle
		$$\bigoplus_{n\in \N} X_n^\bullet \to \bigoplus_{n\in \N} X_n^\bullet \to \colimit_{n\in \N} X_n^\bullet \to \bigoplus_{n\in \N} X_n^\bullet[1]$$
	of the homotopy colimit $\colimit_{n\in \N} X_n^\bullet$. Applying the triangulated functor $\eRHom_A(K^\bullet, -)$ we obtain that
		$$\eRHom_A(K^\bullet, \bigoplus_{n\in \N} X_n^\bullet) \to \eRHom_A(K^\bullet, \bigoplus_{n\in \N} X_n^\bullet) \to \eRHom_A(K^\bullet,\colimit_{n\in \N} X_n^\bullet) \to (\cdots)[1]$$
	is distinguished as well. By enriched compactness of $K$ we have a morphism of distinguished triangles
	$$\begin{tikzcd}[column sep=1em, scale cd=.9]
		\eRHom_A(K^\bullet, \bigoplus_{n\in \N} X_n^\bullet)\ar[r]\ar[d]& \eRHom_A(K^\bullet, \bigoplus_{n\in \N} X_n^\bullet)\ar[r]\ar[d]& \eRHom_A(K^\bullet,\colimit_{n\in \N} X_n^\bullet) \ar[r]\ar[d]& (\cdots)[1]\\
		\bigoplus_{n\in \N} \eRHom_A(K^\bullet, X_n^\bullet)\ar[r]& \bigoplus_{n\in \N} \eRHom_A(K^\bullet,  X_n^\bullet)\ar[r]& \colimit_{n\in \N}\eRHom_A(K^\bullet, X_n^\bullet) \ar[r]& (\cdots)[1]
	\end{tikzcd}$$
	where two out of three vertical morphisms are isomorphism. By \cite[\href{https://stacks.math.columbia.edu/tag/014A}{Tag 014A}]{stacks-project} so is the third and we obtain that
		$$\colimit_{n\in \N}\eRHom_A(K^\bullet, X_n^\bullet) \to \eRHom_A(K^\bullet, \colimit_{n\in \N} X_n^\bullet)$$
	is an isomorphism.
\end{proof}

\section{Compact Generation \& Adjoints}

If a triangulated category is \emph{generated} by derived compact objects then many useful consequences follow. One of these consequences is a very powerful adjoint functor theorem based on Brown representability, originally due to Neeman in \cite{neeman}. Neeman uses this theorem to study Grothendieck duality, essentially providing a proof of the theorem using methods from homotopy theory. We state a slight variation of the adjoint functor theorem in \ref{theorem: an adjoint functor theorem via brown representability} that can be found in \cite[\href{https://stacks.math.columbia.edu/tag/0A8E}{Tag 0A8E}]{stacks-project}. Furthermore, we have to assume that the theorem translates the case where the triangulated category in question is generated by a large(!) class of objects. This is \assumptionPreWord \ref{assumption: another adjoint functor theorem}. 

As in the condensed setting we are often dealing with a large class of compacts we have to be careful in the definition of a compactly generated triangulated category. The standard definition assumes a small collection of generators $(G_i)_{i\in I}$ and requires that the single object $\bigoplus_{i\in I} G_i$ generates the category. The right generalization to a proper class of generators seems to be the analogue of definition \ref{definition: generation}. In the language of \cite[\href{https://stacks.math.columbia.edu/tag/09SI}{Tag 09SI}]{stacks-project} we make the following definition.
\begin{definition}[Compactly generated triangulated categories]
	\label{definition: compactly generated triangulated categories}
	
	Let $\category{D}$ be a triangulated category. A class of $\mathcal G$ of objects of $\category{D}$ is called \emph{generating} if for every object $X\in\category{D}$ there is a small $\mathcal{G'}\subseteq \mathcal G$ such that $X$ is contained in $\langle \mathcal{G'}\rangle$ -- the \emph{smallest strictly full, saturated, triangulated subcategory} of $\category{D}$ containing $\mathcal{G'}$. The triangulated category $\category{D}$ is called \emph{compactly generated} if it admits a class of derived compact generators. We call $\category{D}$ \emph{classically compactly generated} if it can be generated by a small class of derived compact objects.
\end{definition}

Now having defined compactly generated triangulated categories, we can state Brown representability. As already hinted at by the name, this result ensures that certain functors out of the triangulated category in question are representable. This is used in the proof of the adjoint functor theorem to ensure that for each given object one can construct the value of the adjoint functor at that object -- the object being provided by Brown representability. 
\begin{proposition}[Brown representability {\cite[\href{https://stacks.math.columbia.edu/tag/0A8F}{Tag 0A8F}]{stacks-project}}]
	\label{proposition: brown representability}
	
	Let $\category{D}$ be classically compactly generated triangulated category admitting arbitrary direct sums. Let $H\colon \category{D}\to \Ab$ be a contravariant functor which transforms direct sums into products. Then $H$ is representable.
\end{proposition}

This now gives the desired adjoint functor theorem.
\begin{theorem}[An Adjoint Functor Theorem via Brown representability]
	\label{theorem: an adjoint functor theorem via brown representability}
	
	Let $\category{D}$ be a classically compactly generated triangulated category admitting arbitrary direct sums. If $F\colon \category D \to \category D'$ is an exact functor of triangulated categories which preserves direct sums, then $F$ has a right adjoint.
\end{theorem}
\begin{proofsketch}[This is {\cite[\href{https://stacks.math.columbia.edu/tag/0A8G}{Tag 0A8G}]{stacks-project}}]
	
	Let $Y$ be any object of $\category{D}'$ and consider the contravariant functor $\category{D}\to \Ab, X\mapsto \Hom_{\category{D'}}(F(X), Y)$. This functor is cohomological as $F$ is exact and transforms directs sums into products (as $F$ preserves direct sums). Hence, Brown representability \ref{proposition: brown representability} guarantees that this functor is representable, \ie that there exists some object $R(Y)\in\category{D}$ such that $\Hom_\category{D}(-, R(Y)) \cong \Hom_{\category{D'}}(F(-), Y)$. These objects $R(Y)$ then assemble into a functor $R\colon \category{D'}\to \category{D}$.
\end{proofsketch}

The question now is: When does this theorem extend to non-classically compactly generated triangulated categories. Neither is the author aware of such a generalization in the literature nor was he able to prove one. In their notes, Clausen and Scholze always use an(!) adjoint functor theorem \quote{implicitly} when they for example state that \quote{[b]y definition $f_!$ commutes with all direct sums (...) [t]his implies formally that $f_!$ admits a right adjoint $f^!$} at \cite[p. 57]{condenseddotpdf} -- this certainly implies that they have this (or a very similar) adjoint functor theorem in mind.
\begin{assumption}[Another adjoint functor theorem]
	\label{assumption: another adjoint functor theorem}

	Let $\category{D}$ be a compactly generated triangulated category admitting arbitrary direct sums. If $F\colon \category D \to \category D'$ is an exact functor of triangulated categories which preserves direct sums, then $F$ has a right adjoint functor.
\end{assumption}
%\begin{proofsketch}
%	\begin{todo}
%		% TODO
%		This is ew. Make better or remove
%	\end{todo}
%	Denote the class of compact generators of $\category{D}$ by $\mathcal G$. Suppose that $Y$ is any object of $\category{D}'$ and define again the contravariant functor $H \ldef \mleft(X\mapsto \Hom_{\category{D}'}(F(X), Y))\mright)$. Now if $\mathcal G' \subseteq \mathcal G$ is small, denote by $\langle \mathcal G'\rangle$ the smallest full triangulated subcategory of $\category{D}$ with all direct sums, generated by $\mathcal G'$. Then $\langle \mathcal G'\rangle$ is generated by a small set of compacts and $H\vert_{\langle \mathcal G'\rangle}$ still satisfies the conditions for Brown representability, hence is representable. We obtain an object $R_{\langle \mathcal G'\rangle}(Y)\in\category{D}'$ and a natural isomorphism
%	$$\eta_{\langle \mathcal G'\rangle}\colon \Hom_{\category{D'}}(-, R_{\langle \mathcal G'\rangle}(Y))\Rightarrow H\vert_{\langle \mathcal G'\rangle}$$
%	which by Yoneda's lemma corresponds to an element $e_{\langle\mathcal G'\rangle}\in H(R_{\langle\mathcal G'\rangle}(Y))$ which in turns corresponds to a full natural transformation $\Hom_\category{D'}(-, R_{\langle \mathcal G'\rangle}(Y))\Rightarrow H$, again by Yoneda's lemma. This however is not necessarily a natural isomorphism. For any second $\mathcal G'\subseteq \mathcal G'' \subseteq \mathcal G$ small, we obtain similarly
%	$$\eta_{\langle \mathcal G''\rangle}\colon \Hom_{\category{D'}}(-, R_{\langle \mathcal G''\rangle}(Y))\Rightarrow H\vert_{\langle \mathcal G''\rangle}.$$
%	Restricting this second natural isomorphism to $\langle\mathcal G'\rangle$ we obtain a natural isomorphism
%	$$\Hom_{\category{D}'}(-, R_{\langle\mathcal G'\rangle}(Y))\resTo{\langle\mathcal G'\rangle}\xRightarrow{\eta_{\langle\mathcal G'\rangle}} H\resTo{\langle\mathcal G'\rangle} \xRightarrow{\eta_{\langle\mathcal G''\rangle}^{-1}\resTo{\langle\mathcal G'\rangle}} \Hom_\category{D'}(-, R_{\langle\mathcal G''\rangle}(Y))\resTo{\langle\mathcal G'\rangle}.$$
%	Again, by Yoneda this natural transformation corresponds to an element
%	$$f_{\mathcal G', \mathcal G''} \in \Hom_\category{D'}(R_{\langle\mathcal G'\rangle}(Y), R_{\langle\mathcal G''\rangle}(Y))$$
%	whose image under $H$ maps $e_{\langle\mathcal G''\rangle}$ to $e_{\langle\mathcal G'\rangle}$. We obtain a large(!) diagram
%	$$(R_{\langle -\rangle}(Y), f_{-, -})$$
%	over the partially ordered class of small(!) subsets $\mathcal G'\subseteq \mathcal G$. If we suppose that its limit $R(Y)$ exists, then denote by $\tilde e_{\langle\mathcal G\rangle}$ the image of $e_{\langle\mathcal G'\rangle} \in H(R_{\langle\mathcal G'\rangle}(Y))\xrightarrow{\pi_{\mathcal G'}^\ast} H(R(Y))$. Another application of Yoneda yields that the $f_{\mathcal G',\mathcal G''}$ induce natural transformations $H\Rightarrow H$ such that the component $\tilde f_{\mathcal G', \mathcal G''}\colon H(R(Y))\to H(R(Y))$ maps $\tilde e_{\langle\mathcal G'\rangle}$ to $\tilde e_{\langle\mathcal G''\rangle}$. We obtain a second large(!) diagram
%	$$(e_{\langle -\rangle}, \tilde f_{-,-})$$
%	whose colimit determines an element $e\in H(R(Y))$ which in turn by a final application of Yoneda gives a natural transformation
%	$$\Hom_\category{D'}(-, R(Y))\Rightarrow H$$
%	which hopefully is an isomorphism.
%\end{proofsketch}

If one has a pair of functors between two compactly generated triangulated categories, one can say even more. 
\begin{proposition}[{\cite[Thm. 5.1]{neeman}}]
	\label{proposition: preserves compact iff right adjoint preserves coproducts}
	
	Let $F\colon \category{D} \to \category{D'}$ be an exact functor of compactly generated triangulated categories with right adjoint $G$. Then $F$ preserves compact objects if and only if $G$ commutes with direct sums.
\end{proposition}

\begin{remark}
	It seems to the author that the proof of proposition \ref{proposition: preserves compact iff right adjoint preserves coproducts} given by Neeman applies verbatim to the case where $\category{D}$ and $\category{D'}$ are generated by a large class of compacts.
\end{remark}


This implies in particular the following criterion for enriched compactness.
\begin{corollary}[A sufficient condition for enriched compactness]
	\label{corollary: a sufficient condition for enriched compactness}
	
	Let $\category{A}$ be a closed monoidal abelian category that is \grothendieckaxiom{$5$} and is generated by compact projectives. Let $K \in \derived{\category{A}}$ be derived compact. If $-\Dotimes_\category{A} K$ preserves derived compactness then $K$ is enriched compact. This is in particular the case if inside $\category{A}$ the tensor product of any two compact objects is again compact.
\end{corollary}
\begin{proof}
	Observe that $\derived{\category{A}}$ is compactly generated. Since the right adjoint of $-\Dotimes_\category{A} K$ is $\intRHom_\category{A}(K, -)$ by lemma \ref{lemma: deriving a closed monoidal structure} it follows that the right adjoint preserves direct sums (so $K$ is enriched compact) if and only if the left adjoint preserves derived compactness. For the final claim observe that as $\derived{\category{A}}$ is generated by the compact projectives of $\category{A}$ it is enough to check the preservation of derived compactness by $-\Dotimes_\category{A} K$ on these generators since the tensor product commutes with arbitrary direct sums (and hence homotopy colimits). But we can replace $K$ by a complex of such compact projectives and hence a tensor product of the replaced $K$ with a compact generators will be derived compact too.
\end{proof}

\section{A Detour on Cohomology of Topological Spaces}

A sample consequence of the power of derived condensed algebra is that one can recover the sheaf cohomology of any compact Hausdorff space (and hence by extension of any space that is reasonably glued from such compacta). First we need a general definition of cohomology of an object in a (pre)topos.
\begin{definition}[Cohomology internal to a pretopos \& condensed cohomology]
	Let $\category T$ be a pretopos, such that $\Ab(\category T)$ is an abelian category where $\Ext$s exist and such that $\Ab(\category T)\to \category T$ admits a left-adjoint $\Z[-]$. For an object $X\in\category T$ and an abelian group objects $M\in\Ab(\category T)$, define
	$$H_\category T^\bullet(X, M) \ldef \Ext^\bullet_{\Ab(\category T)}(\Z[X], M) \in \derived{\Ab},$$
	the \emph{cohomology of $X$ with coefficients in $M$}. This applies verbatim to condensed sets. Thus, for a condensed set $X$ and a condensed abelian group $M$, define
	$$H_\text{cond}^\bullet(X, M)\ldef H_\condensedSets^\bullet(X, M),$$
	the \emph{condensed cohomology of $X$ with coefficients in $M$}.
\end{definition}

We then have the claimed isomorphism of cohomological theories.
\begin{theorem}[Comparing condensed and sheaf cohomology, {\cite[Thm. 3.2]{condenseddotpdf}}]
	Let $M$ be a discrete topological abelian group. There is an isomorphism
	$$H_\text{sheaf}^\bullet(S, \associated{M}) \ldef \RightD\Gamma(S, \associated{M})\cong \eRHom_{\associated{\Z}}(\freeAbelian{\associated{S}}, \associated{M}) \rdef H_\text{cond}^\bullet(\associated S, \associated M)$$
	natural in $S\in\CHaus$. 
\end{theorem}

%\begin{proof}
%	There is a functor $\alpha^t\colon\open{S}\to \CHaus/S$ such that $U\mapsto (\overline{U} \to S$), giving rise to a map of sites $\alpha\colon \CHaus/S \to \open{S}$. Indeed if $U$ is an open subset of the compact Hausdorff space $S$, then $\overline U$ is compact Hausdorff so the inclusion $\overline U\to S$ is in $\CHaus/S$. Furthermore the functor $\alpha^t$ commutes with finite limits (\ie finite intersections of open sets) and maps covers to covers (it doesnt), since if $U=\bigcup_{i\in I} U_i$ is any cover of $U$
%	
%	Define the functor $\Gamma_\text{cond}(S, -)\ldef \Gamma(S\xrightarrow{\id}S, -)\colon \Sheaves{\CHaus/S}{\Ab} \to \Ab$, to obtain the commuting triangle
%	\begin{center}
%		\begin{tikzcd}
%			\Sheaves{\CHaus/S}{\Ab} &&\Sheaves S \Ab\\
%			&\Ab
%			\ar[from=1-1, to=1-3, "\alpha_\ast"]
%			\ar[from=1-1, to=2-2, "{\Gamma_\text{cond}(S, -)}", swap]
%			\ar[from=1-3, to=2-2, "{\Gamma(S, -)}"]
%		\end{tikzcd}
%	\end{center}
%	where $\alpha_\ast\sheaf F$ is the sheaf mapping $U\in \open{S}$ to $\sheaf F(U\to S)$ and $\Gamma_\text{cond}(S, \sheaf F)=\Gamma(S\to S,\sheaf F)$. Notice how $\alpha_\ast$ is part of a geometric morphism $(\alpha_\ast, \alpha^{-1})$ induced by inclusion of categories $\open{S}\to\CHaus/S$, hence is left-exact. As global sections, also $\Gamma_\text{cond}(S, -)$ and $\Gamma(S, -)$ are left exact. Hence we can form right derived functors and there is a natural morphism
%	$$\RightD \Gamma_\text{cond}(S, -) \longrightarrow \RightD \Gamma(S, -) \circ \RightD \alpha_\ast$$
%	This argument is wrong. Stacks assumes enough injectives. We have none.
%	by the general theory surrounding the Grothendieck spectral sequence \cite[\href{https://stacks.math.columbia.edu/tag/015L}{Tag 015L}]{stacks-project}.	Since $\alpha^{-1}\ladj \alpha_\ast$ is an adjunction between abelian categories, where the left adjoint is exact, the right adjoint preserves injective objects, by \cite[\href{https://stacks.math.columbia.edu/tag/015Z}{Tag 015Z}]{stacks-project}. In particular $\alpha_\ast$ sends injective sheaves to injective (hence $\Gamma(S, -)$-acyclic) sheaves and the above morphism is a natural isomorphism, again by the general theory.
%	
%	Now both $\Sheaves{\CHaus/S}{\Ab}$ and $\Sheaves S \Ab$ admit a constant sheaf $\Z$ by means of the assignments $\Z(X\to S)\ldef\Hom_\Top(X, \Z)$ and $\Z(U)\ldef\Hom_\Top(U, \Z)$ respectively. By definition $\RightD \Gamma(S, \Z)$ gives rise to sheaf cohomology $H_\text{sheaf}^\bullet(S, \Z)$, while $\RightD\Gamma_\text{cond}(S, \Z)$ (is quasi-isomorphic to the complex that) gives rise to the condensed cohomology $H^\bullet_\text{cond}(\associated{S}, \associated{\Z})$ by virtue of the natural isomorphisms
%	$$\Gamma_\text{cond}(S, -)\ldef\Hom_\Top(S, -) \cong \Gamma(S, -) \cong  \Hom_\condensedSets(\associated{S}, -) \cong \Hom_\condensedAb(\Z[\associated{S}], -)$$
%	of abelian groups (where the abelian group structure in each case is pointwise on the domain). If we can show that $(\RightD \alpha_\ast)(\Z) \cong \Z$ in $D(\Ab)$, then
%	$$\RightD \Gamma_\text{cond}(S, \Z) \isorightarrow \RightD \Gamma(S, (\RightD \alpha_\ast)(\Z)) \cong \RightD \Gamma(S, \Z)$$
%	which proves the claim.
%\end{proof}

%\begin{proposition}
%	Let $I$ be any set. Then for any discrete abelian group $M$ there is a natural isomorphism
%	$$\intRHom_{\associated{\Z}}\left(\bigoplus\nolimits_I S^1, M\right) \cong \bigoplus_I M[-1]$$
%	in $\derived{\condensedAb}$. Furthermore $\intRHom_{\associated{\Z}}(\bigoplus_I S^1, \R) = 0$.
%\end{proposition}
%\begin{proof}
%	\begin{todo}
%		Reference this? This needs the whole functorial resolution thing... Maybe just reference that instead and try to prove this.
%	\end{todo}
%\end{proof}

%\begin{lemma}
%	\begin{todo}
%		Proof
%		Difficult! Also: We only need this if we reaaaally want $\intRHom$ and $\Dotimes$ to be derived functors with the first argument fixed. If we ignore that, then we are good to go?
%	\end{todo}
%	Let $\preAnaRing{A}$ be analytic. The category $\mod{\preAnaRing{A}}$ is closed symmetric monoidal with internal $\Hom$ given by $\intHom_\underlying{A}$. Furthermore the compact projective generators $\pfree{A}{S}$  of $\mod{\preAnaRing{A}}$ are internally projective.
%\end{lemma}
%\begin{proof}
%	
%	Let $M$ and $N$ be $\preAnaRing{A}$-complete and let $S$ be extremally disconnected. Then
%	\begin{align*}
%		\Hom_\underlying{A}(\pfree{A}{S}, \intHom_\underlying{A}(X, Y)) &\cong \Hom_\underlying{A}(\pfree{A}{S}\otimes_\preAnaRing{A} X, Y)\\
%		&\cong \Hom_\underlying{A}(X, \intHom_\underlying{A}(\pfree{A}{S}, Y))
%	\end{align*}
%	But $Y$ is $\preAnaRing{A}$-complete, which gives $\intRHom_\underlying{A}(\pfree{A}{S}, Y) \cong \intRHom_\underlying{A}(\ufree{A}{S}, Y)$ giving $\intHom_\underlying{A}(\pfree{A}{S}, Y) \cong \intHom_\underlying{A}(\ufree{A}{S}, Y)$.
%	
%	
%	We need to show exactness of $\intHom_\underlying{A}(\pfree{A}{S}, -) \cong \intHom_\underlying{A}(\ufree{A}{S}, -)$. Thus internal projectivity of $\pfree{A}{S}$ reduces to internal projectivity of $\ufree{A}{S}$ in $\mod{\underlying{A}}$. This holds more generally for any condensed ring $A$. Indeed denote by $S^\delta$ the discrete space with the same underlying space as $S$. Consider the continuous $S^\delta\to S$ with quotient $\ast$. As $\free{A}{-}$ (as a left adjoint) commutes with colimits, the induced sequence
%	$$\bigoplus_S A \to \free{A}{S} \to A \to 0$$
%	is exact after observing that $S^\delta = \coprod_S \{\ast\}$ and $\free{A}{\ast} \cong A$. But $\free{A}{S}$ is projective so the epimorphism $\bigoplus_S A\to \free{A}{S}$ admits a section and hence $\free{A}{S}$ is a retract of $\bigoplus_S A$. Now
%	$$\intHom_A(\bigoplus_S A, -) \cong \prod_S \intHom_A(A, -)\cong \prod_S -$$
%	yields that $\bigoplus_S A$ is internally projective, as direct products are exact. But then $\free{A}{S}$ is a retract of an internally projective module and hence exact itself.
%\end{proof}
